\begin{abstract}
Correcting the luminosity distance of a strongly lensed gravitational-wave (GW) source requires the absolute magnification of its images, which GW data alone cannot measure. We quantify how accurately and how reliably it can be obtained by combining GW signals with an electromagnetic model of the lens galaxy. We simulate $8\times10^5$ lensed binary black holes from the GWTC-4.0 population, four galaxy-lens models from the singular isothermal sphere (SIS) to the elliptical power law with external shear, and an O5-like network and the Einstein Telescope (ET), with GW distance errors calibrated on 167 GWTC events. Using only observables that are unaffected by the mass-sheet degeneracy and contain no distance scale---the image positions relative to the lens light, the GW relative magnification and the image arrival order---estimators trained on the simulations determine the magnification of an ET image to a median 3, 7, 10 and 19 per cent for the four lens models (O5: 10--27 per cent). Adding the time delay and redshifts lowers the errors further, but only because the estimator learns the simulated cosmology. Near caustics the errors grow roughly in proportion to the magnification, and the 90 per cent intervals contain the true value for only about half of the images with $\mu>20$. An SIS lens prior gives confidently wrong intervals for realistic lenses, whereas calibrating on elliptical power-law lenses gives reliable intervals for all lens models at 40 per cent greater width. Bayesian parameter estimation of 120 lensed pairs from simulated detector data confirms the GW error model, and a conventional Bayesian lens reconstruction of the same systems reaches 8 (ET) and 14 (O5) per cent, comparable to the trained estimators; its fast two-step version agrees with the full analysis of the detector data. Positional errors of 0.1\,arcsec for the GW images in the host double the magnification error. All results assume reconstructed image positions of the GW source within the host and no external mass sheet.
\end{abstract}

\begin{keywords}
gravitational lensing: strong -- gravitational waves -- methods: data analysis -- methods: statistical
\end{keywords}
